\section{Introduction}

\paragraph{}
This literature review provides a foundational understanding of image processing in plant phenotyping and directions for developing or investigating solutions for low-cost image-processing phenotyping tools in automated greenhouses. The focus is on the existing image processing solutions applied in low-cost automated greenhouse environments. We begin by examining the fundamental concepts of plant phenotyping and the various traits that are amenable to image processing for predicting plant health, growth, and stress. Specifically, we note that leaf phenotypic traits are a robust proxy indicator of physiological and metabolic factors affecting plant health \cite{zhangHighthroughputPhenotypingPlant2023}.

\paragraph{}
Thereafter, aspects of image processing in plant phenotyping are discussed. Notably, image processing and acquisition, as well as various image sensor technologies and segmentation methods, are reviewed in the context of plant phenotyping. Hyperspectral and RGB imaging are the most common acquisition methods, and segmentation techniques critically influence the overall accuracy, precision, and quality of phenotypical trait analysis \cite{choudhurySegmentationTechniquesChallenges2020, muller-linowPlantScreenMobile2019}.

\paragraph{}
Lastly, the status, challenges, and potential innovations of low-cost automated image processing in greenhouse environments are discussed. In this vein, the literature indicates a lack of affordable and generalizable image processing tools that can analyse and process large amounts of image data on inexpensive, resource-constrained hardware and sensors \cite{gomesComputationalSustainabilityComputational2009, minerviniImageAnalysisNew2015}.